%
% Resumo na língua do documento (obrigatório). Segundo o manual da UNIPAMPA, deve ser escrito em um único parágrafo:
%

\begin{abstract}
A popularização de robôs móveis autônomos de baixo custo depende da integração entre sensores acessíveis, controle embarcado e algoritmos de localização e mapeamento capazes de operar sob restrições reais de processamento, energia e ruído sensorial. Neste contexto, este trabalho tem como objetivo integrar e avaliar uma plataforma robótica física baseada em Raspberry Pi, ROS 2 e câmera RGB-D para execução de SLAM visual em ambiente interno. A pesquisa adota a metodologia \emph{Design Science Research} (DSR), orientando a construção e a avaliação de um artefato tecnológico composto por um robô móvel reaproveitado, interface de baixo nível com microcontrolador, sensores de percepção e uma pilha de software voltada à estimação de movimento, publicação de transformações e geração de mapas. O foco do TCC II é deslocado da modelagem simulada para a validação prática no robô físico, com ênfase na estabilidade do pipeline de odometria, transformações de coordenadas, fusão sensorial e execução do RTAB-Map em hardware limitado. A avaliação proposta considera evidências experimentais como perda de pose, estabilidade do mapa, taxa de processamento, coerência da odometria e impacto de ajustes como o uso de filtro de Kalman estendido para combinar IMU e odometria das rodas. Como resultado parcial, espera-se consolidar uma arquitetura funcional e um protocolo experimental que permitam discutir os limites e possibilidades do SLAM embarcado em uma plataforma robótica de baixo custo, mantendo funcionalidades de navegação, como Nav2, como continuidade futura, e sem tratar simulação ou processamento remoto como componentes centrais da solução.
\end{abstract}

%
% Resumo na outra língua - se o texto for em Português, o abstract é em Inglês e vice-versa -
% como parâmetros devem ser passadas as palavras-chave na outra língua,
% separadas por ponto e vírgula, e grafadas em letras minúsculas, salvo
% em nome próprio, abreviaturas, ..., em que podem ser usadas letras maiúsculas (ver exemplo abaixo).
%

\begin{englishabstract}{visual SLAM; mobile robotics; RGB-D camera; ROS 2; Raspberry Pi}
The adoption of low-cost autonomous mobile robots depends on the integration of accessible sensors, embedded control, and localization and mapping algorithms capable of operating under real constraints of processing, power, and sensor noise. In this context, this work aims to integrate and evaluate a physical robotic platform based on Raspberry Pi, ROS 2, and an RGB-D camera for visual SLAM execution in indoor environments. The research adopts the Design Science Research (DSR) methodology, guiding the construction and evaluation of a technological artifact composed of a reused mobile robot, a low-level microcontroller interface, perception sensors, and a software stack focused on motion estimation, coordinate transform publication, and map generation. The focus of TCC II moves from simulated modeling to practical validation on the physical robot, emphasizing the stability of the odometry pipeline, coordinate transforms, sensor fusion, and RTAB-Map execution on constrained hardware. The proposed evaluation considers experimental evidence such as pose loss, map stability, processing rate, odometry consistency, and the impact of adjustments such as using an Extended Kalman Filter to combine IMU data and wheel odometry. As a partial result, the work is expected to consolidate a functional architecture and an experimental protocol to discuss the limits and possibilities of embedded SLAM on a low-cost robotic platform, keeping navigation features such as Nav2 as future work and not treating simulation or remote processing as central components of the solution.
\end{englishabstract}
